% Terminal apparatus of the concise study: the scope note and the References
% for the sources this companion actually uses. The revision timestamp and the
% rights colophon follow from the shared generation record.

\section{Appendix: Scope and Qualifications}

{\footnotesize
\textbf{Edition, Mass and occurrence.} This concise study treats the Mass of
the First Sunday of Advent, Year~B, as the \work{Roman Missal}, Third Edition,
whose Latin base is the \work{Missale Romanum}, editio typica tertia (2002),
and no.~2 of the \work{Lectionary for Mass for Use in the Dioceses of the
United States of America}, which follows the \work{Ordo lectionum Missae} of
1981, give it for the dioceses of the United States, in the national calendar
alone, for Sunday 29~November 2026. The antiphons, prayers, Preface and solemn
blessing serve the First Sunday of Advent in all three years; the readings are
Year~B's. The solemn blessing is optional. A claim resting on the wording of
an antiphon, the psalm or the Alleluia verse does not hold where another
approved chant replaces it (\work{General Instruction} 48, 61--62, 87).
Nothing here holds for a celebration that replaces this Mass with another.

\textbf{Relation to the full study.} Every claim here is drawn from the full
study of this same Mass, \textit{Rend the heavens and come down: a study of the
proper in three interpretations}, which develops the three interpretations at
length, each with its own four senses and the difficulty on which it strains;
nothing evidence-dependent is added. The four rows on the first page orient a
reader and do not replace those senses.

\textbf{Relations among the texts.} The liturgical books state three
relations: in Advent the readings are chosen for the season, the First
Sunday's Gospel concerning the Lord's coming at the end of time, the
prophet's reading the Messiah and the apostle's an exhortation for the season
(\work{Praenotanda} 67, 93); the psalm responds to the first reading
(\work{General Instruction} 61); and the Alleluia verse greets the Gospel
(\work{General Instruction} 62). The \work{Ordo}'s titles over the three
readings name the plea for God's descent, the waiting for Christ's revelation
and the watch; the \work{Ordo} pairs none of the readings with another. Among
the authors cited, Jerome sets 1~Corinthians 2:9 beside Isaiah 64:3 (Douay
64:4) and Psalm 80 (79):4 beside Isaiah 64:6 (Douay 64:7)
and reads the petition of Isaiah 63:17 also as \latin{converte nos}; Aquinas
sets Psalm 90 (89):13 and Lamentations 5:21 beside the same petition; Augustine
reads Psalm 80 (79) of the Lord's coming; and Augustine, on Mark 13:37, and
Aquinas, on 1~Corinthians 1:8, join each one's death to the last day. Every
other connection between one text of the Mass and another is the editor's
synthesis, grounded where it cites an author in what he says of his own
passage; the agreements of Latin words (\latin{occurrentes} and
\latin{occurristi}; \latin{exspectant} and \latin{exspectare};
\latin{vigilantes} and \latin{vigilate}; \latin{convertere} and \latin{converte
nos}) are observations about the texts. None of the authors cited comments on
this Mass as a whole or links Isaiah 63--64 to Mark~13, and no liturgical
commentary on the Mass is cited.

\textbf{Text, English and rights.} Scripture is quoted in the Douay--Rheims
version as revised by Bishop Challoner, a public-domain study translation of
the Clementine Vulgate, and not in the approved English proclaimed at Mass,
which is that of the Lectionary. Isaiah is cited by the Lectionary's chapters
and verses, which follow the Hebrew, with the Douay--Rheims number beside it
where it differs; psalms are numbered as the Lectionary numbers them, with the
Vulgate number in parentheses. \latin{Descendisti} (Is 64:2b) is the Latin
\work{Ordo}'s; the English proclaimed in the United States, as given on the
bishops' conference's page of readings for this date, has nothing for it. The
Missal's antiphons, prayers, Preface and blessing are identified by Latin
incipit and described; the descriptions are not translations. The Latin and
English of the Missal, the \work{Ordo}, the Lectionary, the \work{General
Instruction} and the \work{Nova Vulgata} are protected: only incipits, the
\work{Ordo}'s titles and response and short Latin phrases needed by the
argument are given. Augustine on the psalms, Chrysostom, Cyril and Bellarmine
are quoted in the public-domain English translations named below; other
English renderings of Latin are the editor's, and the Corpus Thomisticum texts
are quoted in short phrases.

\textbf{Reception and its limits.} Jerome on Isaiah and on Matthew was read on
page images of the Migne printing; Aquinas in the Corpus Thomisticum; Bede and
Gregory in electronic Latin texts of the Migne printing; Augustine's letter 199
in the Latin of augustinus.it and his sermons on Psalm 84 (85) in Latin as well
as English; Bellarmine in O'Sullivan's abridged English, the Latin not
compared. No text was collated with a modern critical edition. Jerome,
Chrysostom and Gregory comment on parallels to the Gospel (Matthew 24, Luke
12), and only Augustine and Bede on Mark 13 itself; Bede's comments on Mark
13:32--33 copy Jerome's on Matthew and are credited to Jerome. Gregory's three
watches belong to Luke's parable and are not Mark's four. Augustine on Psalm
84 (85) is cited by the Latin sections 8, 9 and 15, which the Oxford
translation gives as its §§6 and~12; the sentences of section 9 on Christ given
as man and to be given as God, which that translation omits, are rendered from
the Latin. Jerome's application, near the passages quoted, of Isaiah 63:18 and
64:5 to the Jewish people after the Passion is not used.

\textbf{Chronology and records.} The second page is the generated projection of
the Scripture chronology corpus for this Mass, unchanged; its Isaiah passage is
given in the Clementine numbering. Each passage's sources and setting are in
the full study's appendix Scriptural Date and Location. The texts of the Mass
are audited in \texttt{propers/verified.md}; the search, its loci and its
negative results in \texttt{research/scope.md}; the reasoning behind the three
interpretations in \texttt{research/interpretations.md}; and this companion's
production in \texttt{research/production-review.md}.
\par}

\phantomsection
\section*{References}
\runninghead{References}
\addcontentsline{toc}{section}{References}

{\scriptsize
\subsection*{Liturgical books and Bibles}
\begin{itemize}[itemsep=0pt,topsep=.25em]
\item \work{Missale Romanum}, editio typica tertia (Typis Vaticanis, 2002), \latin{Dominica I Adventus}; \latin{Praefatio I de Adventu}; the solemn blessing \latin{In Adventu}.
\item \work{Ordo lectionum Missae}, editio typica altera (Libreria Editrice Vaticana, 1981): \work{Praenotanda} 67 and 93; no.~2, \latin{Dominica I Adventus}, Year~B.
\item \work{Lectionary for Mass for Use in the Dioceses of the United States of America}, second typical edition, vol.~I, no.~2, as witnessed by the United States Conference of Catholic Bishops' page of readings for 29~November 2026, \url{https://bible.usccb.org/bible/readings/112926.cfm}, read on 26~November 2026.
\item \work{General Instruction of the Roman Missal}, United States edition of 2011 with the emendations of 2021, nos.~48, 61--62 and 87.
\item The Holy Bible, Douay--Rheims version, revised by Bishop Richard Challoner (Project Gutenberg text).
\item \work{Biblia Sacra Vulgatae editionis} (Clementine Vulgate), eBible.org text.
\item \work{Nova Vulgata Bibliorum Sacrorum editio}, as published by the Holy See at vatican.va.
\end{itemize}

\subsection*{Fathers, Doctors and saints}
\begin{itemize}[itemsep=0pt,topsep=.25em]
\item St Jerome, \work{Commentarii in Isaiam} XVII, on 63:16--64:8 (Migne, PL 24, Paris, 1865, cols.~643--650).
\item St Jerome, \work{Commentarii in Matthaeum} IV, on 24:36 and 24:42--43 (Migne, PL 26, Paris, 1845).
\item St Augustine, \work{Enarrationes in Psalmos} 24, 79 and 84, Nicene and Post-Nicene Fathers, 1st ser., vol.~8, and the Latin of the Migne printing for Psalm 84.
\item St Augustine, \work{Epistula} 199, to Hesychius, Latin text of augustinus.it.
\item St John Chrysostom, \work{Homilies on Matthew} 77, Nicene and Post-Nicene Fathers, 1st ser., vol.~10; \work{Homilies on First Corinthians} 2, vol.~12.
\item St Cyril of Jerusalem, \work{Catechetical Lectures} 15, trans. E.~H. Gifford, Nicene and Post-Nicene Fathers, 2nd ser., vol.~7.
\item St Gregory the Great, \work{Homiliae in Evangelia} 9 and 13, Latin text of the Migne printing.
\item St Bede, \work{In Marci evangelium expositio} IV, on 13:32--37, Latin text of the Migne printing.
\item St Thomas Aquinas, \work{Expositio super Isaiam ad litteram}, c.~63--64 (Corpus Thomisticum text).
\item St Thomas Aquinas, \work{Super I epistolam ad Corinthios} c.~1, lect.~1 (Corpus Thomisticum text).
\item St Robert Bellarmine, \work{Explanatio in Psalmos}, as \work{A Commentary on the Book of Psalms}, trans. J. O'Sullivan (1866, abridged), in the eCatholic2000 transcription: on Psalms 79 and 84.
\end{itemize}
\par}
